% Part: first-order-logic
% Chapter: beyond
% Section: other-logics

\documentclass[../../../include/open-logic-section]{subfiles}

\begin{document}

\olfileid[hi]{fol}{byd}{oth}

\olsection{अन्य तर्कशास्त्र}

अब तक आप समझ गए होंगे कि कोई नया तर्कशास्त्र बनाना कठिन नहीं है। आप भी अपना
वाक्यविन्यास बना सकते हैं, एक निगमन-तंत्र गढ़ सकते हैं और उसके अनुरूप
अर्थविज्ञान रच सकते हैं। यदि आप चाहते हैं कि !!{derivation}-तंत्र उस
अर्थविज्ञान के लिए पूर्ण हो, तो आपको कुछ चतुराई दिखानी पड़ सकती है; और
व्यापक जगत को यह मानने के लिए मनाने में भी कुछ श्रम लग सकता है कि आपका
तर्कशास्त्र सचमुच रोचक है। बदले में, आप उसके गणितीय और संगणनात्मक गुणधर्मों
की खोज में घंटों का निर्मल आनंद ले सकते हैं।

हाल के दशकों में औपचारिक तर्कशास्त्रों की मानो बाढ़ आ गई है। अस्पष्ट
गुणधर्मों के विषय में विचार-विन्यास का प्रतिरूपण करने के लिए धुँधला
तर्कशास्त्र बनाया गया है। अनिश्चितता के विषय में विचार-विन्यास के लिए
प्रायिकता-तर्कशास्त्र बनाया गया है। प्रत्याहरणीय विचार-विन्यास---अर्थात्
ऐसे ``उचित'' अनुमानों, जिन्हें नई सूचना मिलने पर बाद में उलटा जा सकता
है---के प्रतिरूपण के लिए पूर्वकल्पना-आधारित तर्कशास्त्र और अ-एकदिशी
तर्कशास्त्र बनाए गए हैं। ज्ञान के विषय में विचार-विन्यास के लिए
ज्ञानमीमांसात्मक तर्कशास्त्र; कार्य-कारण संबंधों के लिए कारणात्मक
तर्कशास्त्र; और नैतिक कर्तव्यों के विषय में विचार-विन्यास के लिए
``कर्तव्यात्मक'' तर्कशास्त्र तक हैं। इन तंत्रों को प्रस्तुत करने की मुख्य
प्रेरणा दार्शनिक, गणितीय अथवा संगणनात्मक है या नहीं, उसके अनुसार आपको ऐसे
प्राणियों का अध्ययन गणितीय तर्कशास्त्र, दार्शनिक तर्कशास्त्र, कृत्रिम
बुद्धिमत्ता, संज्ञान-विज्ञान या किसी अन्य शीर्षक के अंतर्गत मिल सकता है।

सूची बढ़ती ही जाती है और संभावनाएँ अंतहीन लगती हैं। समस्त मानवीय बुद्धि को
परिकलन में घटा देने का लाइबनित्स का स्वप्न हम शायद कभी पूरा न कर पाएँ---पर
यह हमें प्रयास करने से नहीं रोक सकता।

\end{document}
